%
% File: abstract.tex
% Author: Tengxiang Li
%

\chapter{Abstract}
\begin{SingleSpace}
This diploma project presents the development and evaluation of a multimodal emotion recognition system using an attention bottleneck mechanism. The system is designed for the CMU-MOSEI six-emotion multi-label classification task and processes three modalities: text, audio, and visual information. The target emotion categories are happiness, sadness, anger, surprise, disgust, and fear.
The proposed model combines modality-specific encoders with a bottleneck attention fusion module. Textual information is represented using a frozen BERT-base-uncased encoder with averaged hidden-layer representations. Acoustic information is represented using COVAREP features, while visual information is represented using OpenFace facial actionunitfeatures.Themodalityrepresentationsareprojectedintoasharedhiddenspace and fused through learnable bottleneck tokens, which provide an efficient communication channel between modalities without full pairwise cross-modal attention.
The model was trained and evaluated on the CMU-MOSEI dataset using macro averaged F1 score as the main evaluation metric due to class imbalance. Additional valuation was performed using average weighted accuracy and per-class F1 scores. The proposed bottleneck fusion model outperformed the late fusion baseline, achieving an Avg. F1 of 0.4854 with the default threshold and 0.4942 after per-class threshold tuning. The results show that bottleneck-based fusion improves multimodal interaction and provides better recognition of minority emotion classes, especially fear and surprise.
In addition to the experimental model evaluation, the project includes a functional system prototype. A FastAPI backend was implemented to process raw video input by extracting audio, visual, and textual features using ffmpeg, opensmile, py-feat, faster- whisper, and BERT tokenization. The trained model is then used to return emotion probabilities through a REST API. The system demonstrates the practical applicability of the proposed approach as a research prototype for multimodal emotion recognition.

\textbf{Keywords:} multimodal emotion recognition, CMU-MOSEI, attention bottleneck
mechanism, BERT, COVAREP, OpenFace, macro-F1, multi-label classification.

\end{SingleSpace}
% \clearpage
